% !TEX root = ../Projektdokumentation.tex
\section{Entwurfsphase}
\label{sec:Entwurfsphase}

\subsection{Zielplattform}
\label{sec:Zielplattform}
% TODO: kurz beschreiben:
%   - PBS als Zielsystem (Debian-basiert, eigene Web-UI, CLI)
%   - bestehende PVE-Cluster-Umgebung als anzubindendes System


\subsection{Backup-Konzept}
\label{sec:Backupkonzept}
% TODO: hier das ausgearbeitete Backup-Konzept dokumentieren:
%   - Klassifizierung der VMs/CTs (z.\,B. kritisch, normal, Test)
%   - Backup-Frequenzen pro Klasse
%   - Retention Policies pro Klasse
%   - Verschlüsselung / Integritätsprüfung
%   - Restore-Strategie


\subsection{Netzwerk- und Storage-Architektur}
\label{sec:NetzwerkStorage}
% TODO: Architektur beschreiben:
%   - Backup-VLAN und Firewall-Regeln
%   - ZFS-Pool-Layout (RAID-Level, Disks, ARC, Special Devices)
%   - Datastore-Struktur auf dem PBS
% ggfs. ein Architektur-Diagramm im Anhang einfügen.


\subsection{Authentifizierung und Berechtigungen}
\label{sec:AuthZ}
% TODO: Konzept für Authentifizierung beschreiben:
%   - PBS API-Token für PVE
%   - lokale PBS-User vs. LDAP/AD
%   - Berechtigungsmodell (wer darf was)


\subsection{Maßnahmen zur Qualitätssicherung}
\label{sec:QS}
% TODO: Test- und Validierungsstrategie:
%   - Backup-Jobs manuell auslösen
%   - Restore-Tests einzelner VMs/CTs auf Test-Hypervisor
%   - Verifikation der PBS-Datastore-Integrität (verify-jobs)
%   - Monitoring/Benachrichtigung testen
